% Separate proposed strengthening of A.68; no canonical edit.
\paragraph{Every actual intermediate cover of either dihedral closure.}
The same fixed-group conclusion holds for every intermediate
\(D^\circ\to E\setminus\{O\}\) of either degree-sixteen cover.
To prove this, extend the number field so that its actual geometric
subgroup \(H\subset G=D_{16}\), all cusps, and all quotient maps are
defined. Let \(D\) be the smooth proper normalization of the quotient
function field \(\overline K(C)^H\). There are actual finite maps
\[
                   C\xrightarrow{f}D\longrightarrow E,
\]
and the missing points of \(D^\circ\) are precisely the images of
the four cusps of \(C\), equivalently the inverse image of \(O\).
The compact genus of \(D\) is at least one, by Riemann--Hurwitz
for its nonconstant finite map to \(E\).

At \(P_{0,+}\), the elliptic parameter \(-X/Y\) from (A.68.1)
has order four: its leading term is \(2w^4\). The actual inertia
is therefore
\[
 I=\langle r^2\rangle\simeq C_4.
\]
This also follows directly from the displayed deck action:
\(r\) interchanges the two cusps over zero, \(r^2\) fixes
each, and every reflection sends zero to infinity. In particular
\(I\) is normal in \(G\) and is the inertia at every cusp.
Moreover
\[
 \Phi(G)=G^2[G,G]=\langle r^2\rangle=I.
\]
One can check the required Frattini property without invoking
a classification theorem. A subgroup whose image is all of
\(G/I\simeq C_2^2\) contains an odd power of \(r\) and a
reflection: if two independent images are both reflections,
their product gives that odd rotation. These elements generate
\(G\). Consequently \(HI=G\) forces \(H=G\).

The cusps of \(D\) are the actual orbit set
\(H\backslash G/I\). Normality of \(I\) gives its cardinality
\([G:HI]\), which is at least two for every proper subgroup
\(H<G\). For \(H=G\), the intermediate is the original
one-cusp elliptic curve, whose assertion is already proved.

For a proper \(H\), the norm
\(f_*:J_C\to J_D\) is surjective: its composite with pullback
is multiplication by \(\deg f\). Thus \(J_D\) is an isogeny
quotient of \(J_C\). After a finite extension at two it has
good reduction, by the auxiliary-prime
N\'eron--Ogg--Shafarevich criterion, and its special \(2\)-rank
is zero. Indeed an abelian quotient of a good abelian variety
is good, and its rational \(2\)-divisible group is an isogeny
factor; it cannot acquire a slope-zero part absent in \(J_C\).
The proper stable reduction of \(D\) is therefore of compact type.

Let \(P',Q'\) be any two distinct cusps of \(D\). Choose
cusps \(P,Q\) of \(C\) mapping respectively to them. The
relations (A.68.12) give a rational function \(g\) on \(C\)
with \(\operatorname{div}(g)=8(P-Q)\). The function-field
norm satisfies
\[
 \operatorname{div}(\operatorname{Nm}_{K(C)/K(D)}g)
       =f_*\operatorname{div}(g)=8(P'-Q').
\]
There is no ramification multiplier in this pushforward: over
an algebraically closed constant field, the residue degree
of each point is one. This is a norm of an actual function,
not an inference from pullback torsion alone.

Hence every distinct pair of target cusps has difference
annihilated by eight. Its specialization in the good
\(2\)-rank-zero target Jacobian is zero. Separate all marks
in a stable marked model of \(D\); adding rational components
for this separation preserves compact type. The torsion-divisor
criterion of A.41 proves that the unique joining path between
the components carrying \(P'\) and \(Q'\), including both
endpoints, is rational. Since every selected cusp has a
distinct companion, A.42 applies at every coefficient prime.
Its finite-field-restriction argument supplies the indicated
original open-derived subgroups and arbitrary induced packets.
Finally A.59 supplies all normalizers, because the actual
geometric Galois closure of an intermediate is a quotient
of this finite two-group. No pro-odd injectivity of the map
from an intermediate into \(C\) is used.
